\begin{frame}{Machine learning models}

\small

\begin{block}{Machine learning models}

\begin{tabularx}{\linewidth}{@{}p{0.22\linewidth}p{0.32\linewidth}X@{}}

\toprule

\textbf{Model} & \textbf{Description} & \textbf{Scientific question}\\

\midrule

Linear SVR
& Weighted sum of waveform samples
& Is the correction mostly linear?\\[1em]

CNN
& Learns local waveform patterns
& Are nonlinear models better?\\[1em]
\bottomrule

\end{tabularx}

\vspace{0.35em}

\scriptsize
\textit{SVR: Support Vector Regressor \hspace{1.5em} CNN: Convolutional Neural Network}

\end{block}

\vspace{0.3em}

\begin{columns}[T,totalwidth=\textwidth]

\begin{column}{0.26\textwidth}

\begin{block}{Performance}

A model is useful if it reduces CTR on the independent blind set.

\end{block}

\end{column}

\begin{column}{0.66\textwidth}

\begin{alertblock}{Interpretation}

If different models learn similar event-by-event corrections,
performance is likely limited by the information in the signal,
rather than by model complexity.

\end{alertblock}

\end{column}

\end{columns}

\end{frame}